% ─────────────────────────────────────────────────────────────────────────────
\section{Convergents and the Three-Term Recurrence}
% ─────────────────────────────────────────────────────────────────────────────

\subsection{Definition}

\begin{definition}[Convergent]
  The $k$-th \emph{convergent} of $[a_0; a_1, \ldots, a_n]$ is
  $x_k = p_k / q_k$ where the numerator--denominator pairs satisfy the recurrence
  \begin{equation}\label{eq:recurrence}
    \begin{pmatrix} p_k \\ q_k \end{pmatrix}
    = a_k \begin{pmatrix} p_{k-1} \\ q_{k-1} \end{pmatrix}
    + \begin{pmatrix} p_{k-2} \\ q_{k-2} \end{pmatrix}
  \end{equation}
  with initial conditions $p_{-1} = 1$, $p_0 = a_0$, $q_{-1} = 0$, $q_0 = 1$.
\end{definition}

\subsection{Implementation}

\texttt{FinateSimpleContinuedFraction.\_convergent} implements \cref{eq:recurrence}
recursively, re-indexed so that \texttt{body[0]} $= a_1$, \texttt{body[k]} $= a_{k+1}$.
The base cases in the code are:
\begin{align*}
  n=0: &\quad (p_0, q_0) = (1,\; a_1),  \\
  n=1: &\quad (p_1, q_1) = (a_2,\; a_1 a_2 + 1).
\end{align*}

\begin{remark}[Design decision: body-indexed convergents]
The \texttt{convergent(n)} method indexes into \texttt{body}, so $n$ ranges from $0$ to
$\texttt{size}-1$ where $\texttt{size} = |\texttt{body}|$.
The actual rational value requires adding \texttt{head} to $p/q$, which is what
\texttt{convergent\_as\_float} does.
This separation keeps the convergent arithmetic clean and free of a sign-carrying integer part.
\end{remark}

\subsection{Key Properties}

\begin{proposition}[Determinant formula]
  For any $k \geq 0$:
  \[
    p_{k-1} q_k - p_k q_{k-1} = (-1)^k.
  \]
  In particular $\gcd(p_k, q_k) = 1$ for all $k$.
\end{proposition}

\begin{proposition}[Best rational approximation]
  The convergents $p_k/q_k$ are the \emph{best rational approximations} to the
  represented rational in the sense that no fraction with denominator $\leq q_k$
  approximates the value better than $p_k/q_k$.
\end{proposition}

\begin{remark}
  The best-approximation property motivates using SCF convergents for high-precision
  string output: \texttt{convergent\_as\_string} and \texttt{convergent\_float\_as\_string}
  use \texttt{StringMethod.string\_divition}, which computes a long-division decimal
  expansion of a convergent pair $(p_k, q_k)$ to an arbitrary number of digits.
\end{remark}